\section{Discussion}

We presented a three-step, auditable route to Lee--Yang-type critical boundaries for a fixed-order spectral-closure family of partition functions:
\begin{enumerate}
  \item \textbf{Finite-order spectral closure.} A reduced rational generating function yields a minimal order-\(4\) linear recurrence and fixed Hankel rank \(4\) (outside a finite degeneracy set).
  \item \textbf{Discriminant factorization.} The complete factorization of \(\Disc_\lambda\Pi(\lambda,y)\) pins down the algebraic branch locus \(\Branch\), and identifies a unique nontrivial real branch point \(\yLY<0\) that serves as a real-axis edge anchor.
  \item \textbf{Limit log-mgf.} The dominant spectral branch near \(y=1\) controls the limiting log-moment generating function, and implicit differentiation gives closed-form fluctuation constants \(\mu=5/18\) and \(\sigma^2=67/972\) for LLN/CLT.
\end{enumerate}

Several extensions suggest themselves:
\begin{enumerate}
  \item for broader reduced rational families, one may seek a unified, certificate-oriented description of ``discriminant \(\rightarrow\) edge endpoints \(\rightarrow\) equimodular curves'';
  \item under fixed Hankel rank constraints, the computable geometry of root-modulus competition curves merits systematic study;
  \item at the probabilistic level, large deviations and the degeneration of rate functions near edge singularities may provide a finer description of critical fluctuations.
\end{enumerate}

